% !TEX root = task-3-20-1.tex
% Editable task bank. Append new BankTask blocks at the end of the document.
\documentclass[12pt,a4paper]{article}
\usepackage{physics-textbook}
\usepackage{graphicx,needspace}
\usetikzlibrary{decorations.pathmorphing}
\newif\ifShowSolutions
\ifdefined\TutorEdition
  \ShowSolutionstrue
\else
  \ShowSolutionsfalse % Change to \ShowSolutionstrue to display solutions.
\fi
\newcounter{banktask}
\newsavebox{\BankTaskBox}
\newcommand{\BankTaskContent}[3]{%
  \noindent\textbf{\thebanktask.}\quad #2\par
  \ifShowSolutions
    \Needspace{50mm}% Keep the solution heading with its opening text or diagram.
    #3\par
    \ifdefined\TutorEdition
      \ifcsname TutorSource#1\endcsname
        \nopagebreak\smallskip{\small\csname TutorSource#1\endcsname\par}%
      \fi
    \fi
  \fi
}
\newcommand{\BankTask}[3]{%
  \refstepcounter{banktask}\par\medskip
  \sbox{\BankTaskBox}{\begin{minipage}{\linewidth}%
    \BankTaskContent{#1}{#2}{#3}\end{minipage}}%
  \ifdim\dimexpr\ht\BankTaskBox+\dp\BankTaskBox\relax>\dimexpr\textheight-20mm\relax
    \Needspace{35mm}\BankTaskContent{#1}{#2}{#3}%
  \else
    \Needspace{\dimexpr\ht\BankTaskBox+\dp\BankTaskBox\relax}%
    \noindent\usebox{\BankTaskBox}\par
  \fi
  \medskip
}
\begin{document}
\BookSubsection{3.20.1}{Индуктивность}


% Source DB72A1: https://ege.fipi.ru/bank/questions.php?search=1&proj=BA1F39653304A5B041B656915DC36B38&qid=DB72A1
\BankTask{DB72A1}{
Сила тока в катушке с индуктивностью \(0{,}2\,\text{Гн}\) за \(0{,}05\,\text{с}\) равномерно увеличилась от \(2\,\text{А}\) до \(5\,\text{А}\). Чему равен модуль ЭДС самоиндукции в катушке (в $\text{В}$)?
}{
\smallskip\noindent\textbf{Решение от GPT}\par
\[
|\mathcal{E}_{\text{си}}|=L\frac{|\Delta I|}{\Delta t}=0{,}2\cdot\frac{5-2}{0{,}05}=12\,\text{В}.
\]
Ответ: \(12\,\text{В}\).
}

% Source B8FD36: https://ege.fipi.ru/bank/questions.php?search=1&proj=BA1F39653304A5B041B656915DC36B38&qid=B8FD36
\BankTask{B8FD36}{
При равномерном изменении силы тока в катушке на \(10\,\text{А}\) за \(0{,}02\,\text{с}\) в ней возникает ЭДС самоиндукции, равная \(200\,\text{В}\). Чему равна индуктивность катушки (в $\text{Гн}$)?
}{
\smallskip\noindent\textbf{Решение от GPT}\par
Из \(|\mathcal{E}_{\text{си}}|=L|\Delta I|/\Delta t\) получаем
\[
L=\frac{|\mathcal{E}_{\text{си}}|\Delta t}{|\Delta I|}=\frac{200\cdot0{,}02}{10}=0{,}4\,\text{Гн}.
\]
Ответ: \(0{,}4\,\text{Гн}\).
}

% Source 159D34: https://ege.fipi.ru/bank/questions.php?search=1&proj=BA1F39653304A5B041B656915DC36B38&qid=159D34
\BankTask{159D34}{
В катушке индуктивностью \(1\,\text{мГн}\) сила тока в течение \(0{,}1\,\text{с}\) равномерно возрастает от \(0\) до некоторого конечного значения. При этом в катушке наблюдается ЭДС самоиндукции, модуль которой равен \(0{,}2\,\text{В}\). Определите конечное значение силы тока в катушке (в $\text{А}$).
}{
\smallskip\noindent\textbf{Решение от GPT}\par
По закону самоиндукции \(|\mathcal{E}|=L\dfrac{\Delta I}{\Delta t}\). Начальный ток равен нулю, поэтому
\[
I=\frac{|\mathcal{E}|\Delta t}{L}=\frac{0{,}2\cdot0{,}1}{10^{-3}}=20\,\text{А}.
\]
Ответ: \(20\).
}

% Source 31DAF4: https://ege.fipi.ru/bank/questions.php?search=1&proj=BA1F39653304A5B041B656915DC36B38&qid=31DAF4
\BankTask{31DAF4}{
На рисунке приведён график зависимости силы тока от времени в электрической цепи, содержащей катушку, индуктивность которой равна \(1\,\text{мГн}\). Определите модуль ЭДС самоиндукции (в $\text{мкВ}$) в интервале от \(0\) до \(5\,\text{с}\).
\begin{center}
\begin{tikzpicture}[x=0.25cm,y=0.142cm,>=stealth]
\draw[thin] (0,10)--(20,10) (0,20)--(20,20) (0,30)--(20,30);
\draw[thin] (5,0)--(5,30) (10,0)--(10,30) (15,0)--(15,30) (20,0)--(20,30);
\draw[->] (0,0)--(23,0) node[right] {\(t,\,\text{с}\)};
\draw[->] (0,0)--(0,39);
\node[anchor=east] at (-0.8,38) {\(I,\,\text{мА}\)};
\node[below left] at (0,0) {\(0\)};
\foreach \x in {5,10,15,20} \node[below] at (\x,0) {\(\x\)};
\foreach \y in {10,20,30} \node[left] at (0,\y) {\(\y\)};
\draw[thick] (0,0)--(5,30)--(10,20)--(15,20)--(20,0);
\end{tikzpicture}
\end{center}
}{
\smallskip\noindent\textbf{Решение от GPT}\par
В интервале от \(0\) до \(5\,\text{с}\) ток равномерно возрастает от \(0\) до \(30\,\text{мА}\). Тогда
\[
|\mathcal{E}|=L\frac{|\Delta I|}{\Delta t}=10^{-3}\frac{30\cdot10^{-3}}{5}=6\cdot10^{-6}\,\text{В}=6\,\text{мкВ}.
\]
Ответ: \(6\).
}

% Source 9CFC20: https://ege.fipi.ru/bank/questions.php?search=1&proj=BA1F39653304A5B041B656915DC36B38&qid=9CFC20
\BankTask{9CFC20}{
В катушке индуктивностью $20\,\text{мГн}$ сила тока $I$ зависит от времени $t$, как показано на графике, приведённом на рисунке.

\begin{enumerate}
\item Модуль ЭДС самоиндукции в катушке максимален в интервале времени от $0$ до $1\,\text{с}$.
\item Энергия магнитного поля катушки в интервале времени от $2$ до $3\,\text{с}$ равна $40\,\text{мДж}$.
\item Модуль ЭДС самоиндукции в катушке в интервале времени от $3$ до $5\,\text{с}$ равен $10\,\text{мВ}$.
\item Модуль ЭДС самоиндукции в катушке минимален в интервале времени от $3$ до $4\,\text{с}$.
\item Модуль скорости изменения тока в катушке максимален в интервале времени от $5$ до $6\,\text{с}$.
\end{enumerate}

\begin{center}
\begin{tikzpicture}[book diagram,x=0.75cm,y=0.75cm]
\draw[thin] (0,1)--(6,1) (0,2)--(6,2) (0,3)--(6,3) (0,4)--(6,4);
\foreach \x in {1,2,3,4,5,6} {\draw[thin] (\x,0)--(\x,4);}
\draw[-{book tip}] (0,0)--(6.75,0) node[right] {$t,\,\text{с}$};
\draw[-{book tip}] (0,0)--(0,4.9) node[above left] {$I,\,\text{А}$};
\node[below left] at (0,0) {$0$};
\foreach \x in {1,2,3,4,5,6} {\node[below] at (\x,0) {$\x$};}
\foreach \y in {1,2,3,4} {\node[left] at (0,\y) {$\y$};}
\draw[thick] (0,0)--(2,4)--(3,4)--(5,3)--(6,0);
\end{tikzpicture}
\end{center}
}{
\smallskip\noindent\textbf{Решение от GPT}\par
Используем $|\mathcal E|=L|\Delta I/\Delta t|$: на интервалах $0$–$2$, $2$–$3$, $3$–$5$ и $5$–$6\,\text{с}$ модуль скорости изменения тока равен соответственно $2$, $0$, $0{,}5$ и $3\,\text{А}/\text{с}$, поэтому утверждение 1 неверно, а 5 верно. При $2<t<3\,\text{с}$ энергия $W=LI^2/2=0{,}02\cdot4^2/2=0{,}16\,\text{Дж}=160\,\text{мДж}$, поэтому утверждение 2 неверно. На интервале $3$–$5\,\text{с}$ получаем $|\mathcal E|=0{,}02\cdot0{,}5=0{,}01\,\text{В}=10\,\text{мВ}$, поэтому утверждение 3 верно. Минимальная ЭДС равна нулю на интервале $2$–$3\,\text{с}$, поэтому утверждение 4 неверно. Ответ: 3, 5.
}

\end{document}
